% GENERATED FILE. Edit canonical lessons, then run npm run generate:books.
% canonical-source-hash: fnv1a64:4f14d0cb393cca4b
% canonical-lessons: LA-C471-tornus, LA-C471-torno, LA-C471-polio, LA-C471-acuo, LA-C471-lapicida

\chapter{Lathe, To turn on a lathe, To polish, To sharpen, Stonecutter}
\label{ch:la-lathe-to-turn-on-a-lathe-to-polish-to-sharpen-stonecutter}
\hlchaptermodality{471}

\begin{chapteropening}
\textbf{By the end of this chapter:} \emph{I can talk about finishing work in a workshop: the lathe, turning wood, polishing, sharpening tools, and the stonecutter.}

Complete the last lesson of chapter 471: I can talk about finishing work in a workshop: the lathe, turning wood, polishing, sharpening tools, and the stonecutter.
\end{chapteropening}

\section[tornus, tornī]{tornus, tornī — a lathe}

\label{lesson:LA-C471-tornus}



\begin{quote}
Before the new one: say the Latin for a file, then the Latin for to hew.
\end{quote}

\subsection*{You'll want to know: tornus, tornī}
\textbf{tornus, tornī} — \textquotedblleft{}a lathe\textquotedblright{}.

\begin{quote}
Faber in officīnā tornum habet et pōcula lignea facit.
\end{quote}

\begin{quote}
\emph{The craftsman has a lathe in his workshop and makes wooden cups.}
\end{quote}

Second declension, masculine: \emph{tornus, tornī m.} Also the finish of a poem, as if turned on a lathe.

\begin{cousinweb}
\emph{tornus} was borrowed from Greek \emph{tornos} 'lathe, a tool for drawing circles' (EtymDB; Lewis \& Short), from PIE \emph{terh\textsubscript{1}-} 'to rub, to twist'.
\end{cousinweb}

\begin{grammarlens}[title={What you've built}]
One more word for telling, explaining and describing at B1.
\end{grammarlens}

\subsection*{Your turn}
\begin{itemize}
\raggedright
  \item \emph{Say it:} \emph{tornus, tornī}
  \item \emph{Say it:} \emph{tornus, tornī}, once more
  \item \emph{Say it:} a sentence of your own with \emph{tornus, tornī}
  \item \emph{Recall:} say the Latin for a file, then the Latin for to hew, then say \emph{tornus, tornī} again
\end{itemize}

\begin{tcolorbox}[breakable,colback=teal!4,colframe=teal!35!black,title={Before you move on}]
What does tornus, tornī mean? (\textquotedblleft{}A lathe\textquotedblright{}.) Say it once more.
\end{tcolorbox}
\section[tornō, tornāre]{tornō, tornāre — to turn on a lathe, to round off}

\label{lesson:LA-C471-torno}



\begin{quote}
Before the new one: say the Latin for to hew, then the Latin for a lathe.
\end{quote}

\subsection*{You'll want to know: tornō, tornāre}
\textbf{tornō, tornāre} — \textquotedblleft{}to turn on a lathe, to round off\textquotedblright{}.

\begin{quote}
Faber lignum tornat, et pōculum rotundum fit.
\end{quote}

\begin{quote}
\emph{The craftsman turns the wood, and a round cup takes shape.}
\end{quote}

\emph{tornō, tornāre, tornāvī, tornātum}. English \emph{turn} comes from it, borrowed early into Old English.

\begin{cousinweb}
\emph{tornō} is built on \emph{tornus} 'lathe' (EtymDB; Lewis \& Short), a Greek borrowing. English \emph{turn} comes from \emph{tornāre}, borrowed early into Old English and later into French.
\end{cousinweb}

\begin{grammarlens}[title={What you've built}]
One more word for telling, explaining and describing at B1.
\end{grammarlens}

\subsection*{Your turn}
\begin{itemize}
\raggedright
  \item \emph{Say it:} \emph{tornō, tornāre}
  \item \emph{Say it:} \emph{tornō, tornāre}, once more
  \item \emph{Say it:} a sentence of your own with \emph{tornō, tornāre}
  \item \emph{Recall:} say the Latin for to hew, then the Latin for a lathe, then say \emph{tornō, tornāre} again
\end{itemize}

\begin{tcolorbox}[breakable,colback=teal!4,colframe=teal!35!black,title={Before you move on}]
What does tornō, tornāre mean? (\textquotedblleft{}To turn on a lathe, to round off\textquotedblright{}.) Say it once more.
\end{tcolorbox}
\section[poliō, polīre]{poliō, polīre — to polish, to smooth; to refine}

\label{lesson:LA-C471-polio}



\begin{quote}
Before the new one: say the Latin for a lathe, then the Latin for to turn on a lathe.
\end{quote}

\subsection*{You'll want to know: poliō, polīre}
\textbf{poliō, polīre} — \textquotedblleft{}to polish, to smooth; to refine\textquotedblright{}.

\begin{quote}
Servus vāsa argentea polit.
\end{quote}

\begin{quote}
\emph{The slave polishes the silver dishes.}
\end{quote}

\emph{poliō, polīre, polīvī, polītum}. \emph{Polītus} is polished, also refined in manner. English \emph{polish} and \emph{polite} come from it.

\begin{cousinweb}
\emph{poliō} is usually traced to PIE \emph{pelh\textsubscript{2}-} (EtymDB), though Lewis \& Short suggest another root and the origin is not settled. English \emph{polish} comes from it through French, and \emph{polite} from its participle \emph{polītus}.
\end{cousinweb}

\begin{grammarlens}[title={What you've built}]
One more word for telling, explaining and describing at B1.
\end{grammarlens}

\subsection*{Your turn}
\begin{itemize}
\raggedright
  \item \emph{Say it:} \emph{poliō, polīre}
  \item \emph{Say it:} \emph{poliō, polīre}, once more
  \item \emph{Say it:} a sentence of your own with \emph{poliō, polīre}
  \item \emph{Recall:} say the Latin for a lathe, then the Latin for to turn on a lathe, then say \emph{poliō, polīre} again
\end{itemize}

\begin{tcolorbox}[breakable,colback=teal!4,colframe=teal!35!black,title={Before you move on}]
What does poliō, polīre mean? (\textquotedblleft{}To polish, to smooth; to refine\textquotedblright{}.) Say it once more.
\end{tcolorbox}
\section[acuō, acuere]{acuō, acuere — to sharpen, to whet}

\label{lesson:LA-C471-acuo}



\begin{quote}
Before the new one: say the Latin for to turn on a lathe, then the Latin for to polish.
\end{quote}

\subsection*{You'll want to know: acuō, acuere}
\textbf{acuō, acuere} — \textquotedblleft{}to sharpen, to whet\textquotedblright{}.

\begin{quote}
Mīles gladium ante bellum acuit.
\end{quote}

\begin{quote}
\emph{The soldier sharpens his sword before the war.}
\end{quote}

\emph{acuō, acuere, acuī, acūtum}. Its participle \emph{acūtus} 'sharp' gives English \emph{acute}.

\begin{cousinweb}
\emph{acuō} is inherited from PIE \emph{h\textsubscript{2}e\'{k}-} 'sharp' (EtymDB), the root of \emph{acus} 'needle'. English \emph{acute} comes from its participle \emph{acūtus}, and \emph{acumen} from a noun built on it.
\end{cousinweb}

\begin{grammarlens}[title={What you've built}]
One more word for telling, explaining and describing at B1.
\end{grammarlens}

\subsection*{Your turn}
\begin{itemize}
\raggedright
  \item \emph{Say it:} \emph{acuō, acuere}
  \item \emph{Say it:} \emph{acuō, acuere}, once more
  \item \emph{Say it:} a sentence of your own with \emph{acuō, acuere}
  \item \emph{Recall:} say the Latin for to turn on a lathe, then the Latin for to polish, then say \emph{acuō, acuere} again
\end{itemize}

\begin{tcolorbox}[breakable,colback=teal!4,colframe=teal!35!black,title={Before you move on}]
What does acuō, acuere mean? (\textquotedblleft{}To sharpen, to whet\textquotedblright{}.) Say it once more.
\end{tcolorbox}
\section[lapicīda, lapicīdae]{lapicīda, lapicīdae — a stonecutter, a quarryman}

\label{lesson:LA-C471-lapicida}



\begin{quote}
Before the new one: say the Latin for to polish, then the Latin for to sharpen.
\end{quote}

\subsection*{You'll want to know: lapicīda, lapicīdae}
\textbf{lapicīda, lapicīdae} — \textquotedblleft{}a stonecutter, a quarryman\textquotedblright{}.

\begin{quote}
Lapicīda saxum magnum in forō caedit.
\end{quote}

\begin{quote}
\emph{The stonecutter cuts a large stone in the forum.}
\end{quote}

First declension, masculine: \emph{lapicīda, lapicīdae m.} Like \emph{agricola}, a masculine noun in \emph{-a}.

\begin{cousinweb}
\emph{lapicīda} is built from \emph{lapis, lapidis} 'stone' and \emph{caedō} 'to cut' (Lewis \& Short; an older form is \emph{lapidicīda}). EtymDB lists it without a parent.
\end{cousinweb}

\begin{grammarlens}[title={What you've built}]
One more word for telling, explaining and describing at B1.
\end{grammarlens}

\subsection*{Your turn}
\begin{itemize}
\raggedright
  \item \emph{Say it:} \emph{lapicīda, lapicīdae}
  \item \emph{Say it:} \emph{lapicīda, lapicīdae}, once more
  \item \emph{Say it:} a sentence of your own with \emph{lapicīda, lapicīdae}
  \item \emph{Recall:} say the Latin for to polish, then the Latin for to sharpen, then say \emph{lapicīda, lapicīdae} again
\end{itemize}

\begin{tcolorbox}[breakable,colback=teal!4,colframe=teal!35!black,title={Before you move on}]
What does lapicīda, lapicīdae mean? (\textquotedblleft{}A stonecutter, a quarryman\textquotedblright{}.) Say it once more.
\end{tcolorbox}
